% Related Work: organised by theme. Each subsection ends by stating
% what this paper does differently.

\paragraph{Hierarchical and self-editing agent memory.}
A family of systems extends an agent's effective context by staging
state outside the window. MemGPT \citep{packer2023memgpt} frames this
as virtual memory: the model issues explicit swap-in and swap-out
calls between a main context and a paged store. A-MEM
\citep{xu2025amem} organises memories as a dynamic graph of linked
notes in the style of Zettelkasten, letting the agent grow its own
index. MemoryBank \citep{zhong2024memorybank} introduces an
Ebbinghaus-inspired forgetting curve, and Mem0
\citep{chhikara2025mem0} productionises a scalable long-term memory
layer for deployed agents. RecurrentGPT \citep{zhou2023recurrentgpt}
separates short- and long-term natural-language memory for sustained
generation. All of these expose memory through a retrieval or paging
interface, with addresses chosen by a controller rather than by the
agent. Our approach differs in that the address space is
\emph{user-visible and path-structured}: the agent reasons over paths
the same way a programmer reasons over directories, and chooses
disclosure depth per entry rather than receiving a fixed slice.

\paragraph{Reasoning-and-acting loops and cognitive architectures.}
Agents need context, but they also need a protocol that uses it.
ReAct \citep{yao2023react} interleaves reasoning and action in a
single trace; Reflexion \citep{shinn2023reflexion} adds verbal
self-critique as an episodic memory; Tree of Thoughts
\citep{yao2023tot} generalises the trace into a search tree; and
Generative Agents \citep{park2023generative} combine a memory stream
with periodic reflection. CoALA \citep{sumers2024coala} frames these
designs as cognitive architectures with explicit working and
long-term memory, and Voyager \citep{wang2023voyager} shows that an
ever-growing skill library can serve as durable context for an
embodied agent. Chain-of-thought prompting \citep{wei2022cot}
established that the content of the prompt drives reasoning quality.
Our context layer is complementary: it supplies a substrate that these
loops can read, write, and audit without imposing a particular control
policy.

\paragraph{Long-context models and their limits.}
A tempting alternative to structured memory is simply a longer
window. \citet{liu2024lostmiddle} show that recall is worse for items
placed in the middle of long contexts, and LongBench
\citep{bai2024longbench} quantifies how benchmark accuracy trails the
nominal window size. StreamingLLM \citep{xiao2024streamingllm}
maintains attention sinks for indefinite inference but does not
address \emph{what} to place in the kept window. These findings
justify structured context management even when windows are large:
knowing what to include is a distinct problem from being able to
include it.

\paragraph{Retrieval, tools, and agent frameworks.}
Retrieval-augmented generation \citep{lewis2020rag} injects retrieved
passages into the prompt; Self-RAG \citep{asai2024selfrag} makes the
retrieval decision adaptive via reflection tokens; and HippoRAG
\citep{gutierrez2024hipporag} couples retrieval with a graph of
entities to approximate long-term memory. Toolformer
\citep{schick2023toolformer} showed that models can learn to call
external tools, and multi-agent frameworks such as AutoGen
\citep{wu2023autogen} orchestrate agents that share state through
conversational channels. In these systems, context management is
typically an unstated concern: documents arrive, tools are listed,
conversation history accumulates. Our work promotes context
management itself to a first-class subsystem, with a defined address
space, disclosure levels, and audit log.

\paragraph{Positioning.} The closest prior work is MemGPT
\citep{packer2023memgpt} and A-MEM \citep{xu2025amem}. MemGPT gives
the agent a main/archival split and paging calls; it does not expose
\emph{where} a memory lives to the agent. A-MEM gives a graph, but
links are learned and addresses are opaque. Our contribution is to
combine a transparent path-based address space with per-entry
multi-level disclosure, so that navigation and detail selection are
both explicit, inspectable decisions --- and to back this with an
event-sourced log that lets us reproduce and evaluate context use
turn by turn.
